\section{Computable Sequential PAC-Bayes Certificate}
\label{sec:certificate}

\subsection{Augmented Hierarchical Prior}
Condition on the information $\mathcal G_0$ available after $\mathcal D_{\mathrm{prior}}$. The hypothesis space is augmented with the discrete labels used to construct the Gaussian component:
\begin{equation}
 h=(M,v,s,a),\qquad M=(F,p,\gamma,r,K),
\label{eq:augmented_hypothesis}
\end{equation}
where $v$ denotes the prior-localization variant, $s$ the prior-scale component, and $a$ the consequent vector. The prediction depends on $(M,a)$; the labels $(v,s)$ make the mixture accounting explicit.

Let $\pi_F$, $\pi_p$, and $\pi_\gamma$ be uniform on the predeclared family, lag, and ridge-penalty grids. For ridge and fixed-$K$ families, the radius factor is one; for radius-controlled families, $\pi_{r\mid F}$ is uniform on the radius grid. TSK structures additionally receive the normalized geometric rule-count mass
\begin{equation}
\pi_K(K)=\frac{\exp(-\eta K)}{\sum_{j=1}^{K_{\max}}\exp(-\eta j)},\qquad \eta=\log 2.
\label{eq:kprior}
\end{equation}
The structural mass used by the implementation is therefore
\begin{equation}
\pi(M)=\pi_F(F)\pi_p(p)\pi_\gamma(\gamma)\pi_{r\mid F}(r\mid F)\,\pi_{K\mid F}(K\mid F),
\label{eq:structural_prior}
\end{equation}
where $\pi_{K\mid F}=1$ for ridge and $\pi_{K\mid F}=\pi_K$ for TSK families. For a radius candidate, $K=K(r,\mathcal D_{\mathrm{prior}})$ is fixed after conditioning on $\mathcal G_0$. Ineligible radius candidates receive zero executable mass. For each radius-controlled family,
\begin{equation}
\sum_{r}\pi_{r\mid F}(r\mid F)\pi_K\!\left(K(r)\right)\mathbf 1\{r\text{ eligible}\}\le 1,
\end{equation}
and for any evaluated fixed-$K$ subset, $\sum_{K\in\mathcal K_{\mathrm{exec}}}\pi_K(K)\le1$. The unused family mass is assigned to a dummy hypothesis. Thus the prior is proper without renormalizing after observing which candidates are executable; the reported code lengths are conservative.

For localization variant $v$ and prior scale $s$, let $\omega(v)$ and $\xi(s)$ denote the predeclared mixture masses. The operational studies use a single chronological localized variant, while the development localization ablation uses two equally weighted variants. The scale set is $\mathcal S=\{0.5,1,2,4,8\}$ with uniform mass $\xi(s)=1/5$. For every $(M,v,s)$, define
\begin{equation}
P_{M,v,s}=\mathcal N\!\left(\mu_{P,M,v},\operatorname{diag}(s^2b_{M,v}^2)\right).
\label{eq:prior_gauss}
\end{equation}
The full prior on the augmented space is
\begin{equation}
P(dM,dv,ds,da)=\pi(M)\omega(v)\xi(s)P_{M,v,s}(da),
\label{eq:joint_prior}
\end{equation}
with the remaining structural mass placed on the dummy component.

After observing the certification data, a selected posterior has the form
\begin{equation}
Q=\delta_{\widehat M}\delta_{\widehat v}\delta_{\widehat s}\,
Q_{\widehat M,\widehat v,\widehat s,\widehat\rho},
\label{eq:joint_posterior}
\end{equation}
where
\begin{equation}
Q_{M,v,s,\rho}=\mathcal N\!\left(\mu_{Q,M},\operatorname{diag}(s^2\rho^2b_{M,v}^2)\right),
\label{eq:posterior_gauss}
\end{equation}
and $\rho\in\{0.01,0.03,0.1,0.3,1\}$. Because the component labels are explicit coordinates of the augmented hypothesis, the divergence decomposes exactly as
\begin{align}
\mathrm{KL}(Q\|P)
={}&-\log\pi(\widehat M)-\log\omega(\widehat v)-\log\xi(\widehat s)\nonumber\\
&+\mathrm{KL}\!\left(Q_{\widehat M,\widehat v,\widehat s,\widehat\rho}\middle\|P_{\widehat M,\widehat v,\widehat s}\right).
\label{eq:joint_kl}
\end{align}
The final term is the exact diagonal-Gaussian divergence
\begin{equation}
\frac12\sum_{j=1}^{D_M}\left[
\log\frac{\sigma_{P,j}^2}{\sigma_{Q,j}^2}-1+
\frac{\sigma_{Q,j}^2+(\mu_{Q,j}-\mu_{P,j})^2}{\sigma_{P,j}^2}
\right].
\label{eq:gaussian_kl}
\end{equation}
Selection of $\rho$ needs no additional discrete union penalty because Proposition~\ref{prop:bound} is simultaneous over all $Q\ll P$. Temperature selection is different: it is a choice among concentration inequalities and is charged through the temperature mass $\nu(\lambda)$. The supplementary material gives the proper-mass completion and the componentwise derivation in full.

\subsection{Time-Uniform Bounded-Loss Specialization}
For $h=(M,a)$, write $\mu_t(h)=\mathbb E[\ell_t(h)\mid\mathcal G_{t-1}]$ and
\begin{align}
\widehat R_t(Q)&=\frac1t\sum_{i=1}^{t}\mathbb E_Q[\ell_i],\\
R_t^{\mathrm{seq}}(Q)&=\frac1t\sum_{i=1}^{t}\mathbb E_Q[\mu_i].
\label{eq:risks}
\end{align}

\begin{assumption}[Predictable bounded loss]
For each fixed $h$, the forecast at time $t$ is $\mathcal G_{t-1}$-measurable and $0\le\ell_t(h)\le1$ almost surely.
\end{assumption}

\begin{proposition}[Bounded-loss time-uniform specialization]
\label{prop:bound}
Let $P$ be $\mathcal G_0$-measurable and let $\Lambda$ be a finite positive temperature grid with mass function $\nu$. Conditionally on $\mathcal G_0$, with probability at least $1-\delta$, simultaneously for every $t\ge1$, every $\mathcal G_t$-measurable $Q_t\ll P$, and every $\lambda\in\Lambda$,
\begin{equation}
R_t^{\mathrm{seq}}(Q_t)\le \widehat R_t(Q_t)+\frac{\lambda}{8}+
\frac{\mathrm{KL}(Q_t\|P)+\log[1/(\delta\nu(\lambda))]}{\lambda t}.
\label{eq:main_bound}
\end{equation}
\end{proposition}

Proposition~\ref{prop:bound} is the conditional-Hoeffding instance of the established PAC-Bayes supermartingale recipe \cite{seldin2012martingales,chugg2023unified}. For fixed $h$, $\exp\{\lambda\sum_{i\le t}(\mu_i-\ell_i)-t\lambda^2/8\}$ is a nonnegative supermartingale. Mixing over $P$, applying Ville's inequality, then the Donsker--Varadhan change of measure proves the statement; a union bound charges the finite temperature grid. The proof and the prior-mixture completion are given in the supplementary material.

The reported value is the minimum untruncated right-hand side over the fixed temperature grid, clipped at one only for presentation. Familywise studies allocate $\delta/S$ to each of $S$ predeclared series.

\subsection{Analytic Empirical-Risk Upper Bound}
For fixed antecedents, $f_{M,a}(x)=\phi_M(x)^\top a$ before projection. Therefore
\begin{equation}
\mathbb E_Q[(\bar Y-\phi^\top a)^2]=(\bar Y-\phi^\top\mu_Q)^2+
\phi^\top\operatorname{diag}(\sigma_Q^2)\phi.
\label{eq:secondmoment}
\end{equation}
Projection cannot increase squared distance to $\bar Y\in[-B,B]$. Using the concavity of $u\mapsto\min(1,u)$ gives the pointwise upper bound
\begin{equation}
\widehat r_i(Q)=\min\left\{1,
\frac{(\bar Y_i-\phi_i^\top\mu_Q)^2+
\phi_i^\top\operatorname{diag}(\sigma_Q^2)\phi_i}{4B^2}
\right\}.
\label{eq:empirical_upper}
\end{equation}
The implementation substitutes $t^{-1}\sum_i\widehat r_i(Q)$ for $\widehat R_t(Q)$, which is conservative.

\begin{corollary}[Clipped Bayesian model average]\label{cor:bma}
The clipped Bayesian model-average risk is no larger than the Gibbs risk in Proposition~\ref{prop:bound}, by convexity of squared loss.
\end{corollary}

\subsection{Scope of the Guarantee}
The guarantee is for average conditional risk of the clipped Gibbs predictor, and by convexity its clipped model average, over the observed certification horizon. It is valid for adapted overlapping windows and does not require an estimated mixing coefficient. It is not a guarantee for raw RMSE, an arbitrary deterministic parameter-mean implementation, the future test horizon, or performance relative to a seasonal baseline. Those quantities remain empirical evaluations and are deliberately separated from the certificate.
